\clearpage\thispagestyle{empty}
\pdfbookmark[0]{Sources for Gauss to Bode (9 August 1815)}{bode1815aug-sources}
\begin{center}\textbf{Sources for Gauss to Bode (9 August 1815)}\end{center}
\noindent This AI translation follows Carl Friedrich Gauss’s letter to Johann Elert Bode in \href{https://gdz.sub.uni-goettingen.de/id/PPN236007467}{\textit{Werke}, vol.~VI} (Göttingen, 1874), pp.~387--389. Bode printed it with a heading on the comet observations, orbital elements and observations of Juno in the \textit{Astronomisches Jahrbuch für 1818} (Berlin, 1815), pp.~\href{https://api.digitale-sammlungen.de/iiif/image/v2/bsb10538347_00237/full/1600,/0/default.jpg}{229}--\href{https://api.digitale-sammlungen.de/iiif/image/v2/bsb10538347_00240/full/1600,/0/default.jpg}{232}. The links lead to the Bavarian State Library copy.
\par\smallskip\noindent The yearbook gives 59~13~3 for the comet’s right ascension on 30~March (p.~\href{https://api.digitale-sammlungen.de/iiif/image/v2/bsb10538347_00238/full/1600,/0/default.jpg}{230}), while the collected letter gives 59~13~2. On p.~\href{https://api.digitale-sammlungen.de/iiif/image/v2/bsb10538347_00239/full/1600,/0/default.jpg}{231} the yearbook has \textit{Elemente}, where the collected edition prints \textit{Elemnnte}, rendered here as “elements.” Both printings give 20~37 followed by a colon in the declination for 27~July.
\par\smallskip\noindent On p.~\href{https://api.digitale-sammlungen.de/iiif/image/v2/bsb10538347_00240/full/1600,/0/default.jpg}{232} the yearbook gives the period as $77\tfrac13$ years and Juno’s declination on 29~March as $1\,32\,30{,}1$ N. The collected letter gives $77\tfrac12$ years and $1\,32\,38.1$ N. In the elliptical elements the yearbook says “distance at perihelion” and gives the ascending node without unit marks. The collected edition says “distance from perihelion” and marks the node as $83^\circ\,26'\,21''$. Both printings give half an arcminute for the mount. The spelling \textit{Frauenhofer} follows the collected edition. Translation and typesetting produced by AI.
\clearpage
